\subsection{The \(\zeta(3)\) and \(\zeta(4)\) construction}

We write $(\mathcal M,x_0,v)_F$ for the forward scheme and
$(\mathcal M,x_0,v)_D$ for the dual scheme of a CMF $\mathcal M$ with
starting point $x_0$ and direction $v$.

\begin{theorem}
\label{thm:z34-limit}
Let $A,B\in\mathbb Z$, $x_0=(0,0,1,1,1;1,1,1,1)$ and
$v_F=(-14,-15,13,14,12;1,2,3,4)$.
The forward scheme
\[
\bigl({}_5F_4(1),\,x_0,\,v_F\bigr)_F
\]
has initial condition and projection vector
\[
I_F=\begin{pmatrix}0&0&-6A&3A-2B\\12&0&0&0\end{pmatrix},
\qquad v=\begin{pmatrix}1\\0\\0\\0\end{pmatrix}.
\]
Writing $W_F(N)=F_{x_0,v_F}(0)\cdots F_{x_0,v_F}(N-1)$, the ratio
\[
R_N=\frac{(I_FW_F(N)v)_1}{(I_FW_F(N)v)_2}
\]
converges to $A\zeta(3)+B\zeta(4)$, with exponential error bound
\[
\limsup_{N\to\infty}|R_N-A\zeta(3)-B\zeta(4)|^{1/N}
\leq\tau:=\frac{|\rho_4|}{|\rho_1|}
\approx1.019\cdot10^{-49},
\]
where the roots are those of Proposition~\ref{prop:z34-rates}.
\end{theorem}
Unlike the preceding case, we cannot identify the full dual Meijer G
scheme with this forward scheme by an invertible rational coboundary.
Instead, a rank-three intertwining matrix connects the selected walks.
We show that this connection preserves the dominant dual contribution
and the initial-state identity determining the limit.

\begin{proposition}
\label{prop:z34-coboundary}
Let $y_0=(0,0,1,1,2;0,0,0,0,0)$ and
$v_G=(-16,-17,11,12,10;-2,-1,0,1,2)$.
The dual scheme $(G^{5,2}_{5,5}(1),y_0,v_G)_D$, restricted to its first
three coordinates, is coboundary equivalent, up to scalar normalization,
to the forward scheme $({}_5F_4(1),x_0,v_F)_F$ on a three-dimensional
invariant space $\operatorname{im}E(n)$.
\end{proposition}

\begin{proof}
The coboundary $H(n)$, basis matrix $E(n)$ for the invariant subspaces,
and scalar factors $A(n)$ and $B(n)$ are given in the supplied code,
which verifies
\[
F_{x_0,v_F}(n)E(n+1)=E(n)\widehat F(n),
\qquad
\widehat F(n)=\frac{A(n)}{B(n)}H(n)Q(n)^{-T}H(n+1)^{-1}.
\]
Here $Q(n)$ is the leading $3\times3$ block of $G_{y_0,v_G}(n)$,
$\widehat F(n)=LF_{x_0,v_F}(n)E(n+1)$, and $L$ selects the first, third
and fourth coordinates, with $LE(n)=I_3$.
The restriction of the dual Meijer G matrix follows from
\[
G_{y_0,v_G}(n)=\begin{pmatrix}Q(n)&b(n)\\0&-1\end{pmatrix},
\qquad
G_{y_0,v_G}(n)^{-T}
=\begin{pmatrix}Q(n)^{-T}&0\\b(n)^TQ(n)^{-T}&-1\end{pmatrix}.
\]
Thus a row $(u,0)$ is mapped to $(uQ(n)^{-T},0)$, so the first
three coordinates form an invariant space for the dual walk.
Put
\[
J=\begin{pmatrix}1&0&0&0\\0&1&0&0\\0&0&1&0\end{pmatrix},
\qquad K(n)=E(n)H(n)J,
\qquad R(n)=\frac{A(n)}{B(n)}.
\]
The matrix $K(n)$ has rank three, and the preceding identities imply
\[
F_{x_0,v_F}(n)K(n+1)=R(n)K(n)G_{y_0,v_G}(n)^{-T}.
\]
Consequently, writing $W_G(N)=G_{y_0,v_G}(0)\cdots
G_{y_0,v_G}(N-1)$, we obtain the exact walk identity
\[
W_F(N)K(N)
=\left(\prod_{n=0}^{N-1}R(n)\right)K(0)W_G(N)^{-T}.
\]
This identity does not require $K(n)$ to be invertible.
\end{proof}
Along the trajectory define
\[
\mathcal{G}_n(z)=G_{5,5}^{5,2}\!\left(
 z\;\middle|\;
 \begin{matrix}
 -16n,-17n;\ 1+11n,1+12n,2+10n\\
 -2n,-n,0,n,2n
 \end{matrix}\right),
\qquad \theta=z\,\frac{d}{dz}.
\]

\begin{proposition}
\label{prop:z34-initial-values}
Let
\[
\mathcal{S}_G(n)=\bigl(
\mathcal{G}_n(1),\ (\theta\mathcal{G}_n)(1),\
(\theta^2\mathcal{G}_n)(1),\ (\theta^3\mathcal{G}_n)(1)\bigr),
\]
\[
I_G=\begin{pmatrix}2B&-3A-2B&3A&0\\0&0&6&0\end{pmatrix}.
\]
Then
\[
I_G\mathcal{S}_G(0)^T
=6\begin{pmatrix}A\zeta(3)+B\zeta(4)\\1\end{pmatrix}.
\]
\end{proposition}

\begin{proof}
The Mellin--Barnes integral in the proof of
Proposition~\ref{prop:z34-growth}, with $n=0$, and the reflection formula
\cite[Eq.~5.5.3]{NIST:DLMF} give
\[
\mathcal{G}_0(z)=\frac{1}{2\pi i}\int_{\frac12+i\mathbb R}
\frac{\pi^2z^{-s}}{\sin^2(\pi s)s^3(s+1)}\,ds.
\]
Applying $\theta^k$ multiplies the integrand by $(-s)^k$.
At $z=1$, close the contour to the right along lines
$\operatorname{Re}s=m+\tfrac12$. For $0\leq k\leq3$ their integrals are
$O(m^{k-4})\to0$; the horizontal integrals vanish exponentially.
The poles at positive integers are double, and the clockwise orientation gives
\[
(\theta^k\mathcal{G}_0)(1)
=(-1)^{k+1}\sum_{\ell=1}^{\infty}
\left[\frac{d}{ds}\frac{1}{s^{3-k}(s+1)}\right]_{s=\ell}.
\]
Partial fractions give
\[
\mathcal{S}_G(0)=\bigl(
1-2\zeta(3)+3\zeta(4),\ 1-2\zeta(3),\ 1,\ 1-\zeta(2)\bigr),
\]
and multiplication by $I_G$ gives the claimed identity.
\end{proof}

\begin{proposition}
\label{prop:z34-rates}
The trajectory matrix $G_{y_0,v_G}(n)$ has exponential growth rates given by
two conjugate dominant roots
\[
\rho_{1,2}\approx(2.803\pm6.487i)\cdot10^{25},
\qquad \rho_3=-1,
\qquad \rho_4\approx7.203\cdot10^{-24}.
\]
No factorial normalization is needed. The characteristic roots at infinity
of the restricted matrix $Q(n)$ are $\rho_1,\rho_2,\rho_4$.
\end{proposition}

\begin{proof}
The roots are those of the characteristic polynomial at infinity, extracted
in the supplied code. The factor corresponding to $\rho_3=-1$ is removed
by the block restriction in Proposition~\ref{prop:z34-coboundary}.
\end{proof}

\begin{proposition}
\label{prop:z34-growth}
Let
\[
\eta=\frac{16^{16}15^{30}2^2 3^3 4^4}{13^{13}12^{12}30^{30}}
\left(\frac32\right)^{55}\approx8.517\cdot10^{-4}<1.
\]
Then
\[
(\theta^k\mathcal{G}_n)(1)=O(n^k\eta^n),
\qquad k=0,1,2,3.
\]
\end{proposition}

\begin{proof}
Put $L=14n+1$ and $C_n=(14n)!(16n)!/[(11n)!^2(8n+1)!]$.
The Mellin--Barnes representation \cite[Eq.~16.17.1]{NIST:DLMF} is
\[
\mathcal{G}_n(z)=\frac{1}{2\pi i}
\int_{2n+\frac12-i\infty}^{2n+\frac12+i\infty}
\frac{
 \Gamma(16n+1-s)\Gamma(17n+1-s)
}{
 \Gamma(s+11n+1)\Gamma(s+12n+1)\Gamma(s+10n+2)
}\prod_{j=-2}^{2}\Gamma(s+jn)\,z^{-s}\,ds.
\]
The contour is traversed from bottom to top. Its left pole families
lie in $s\leq2n$ and its right pole families in $s\geq16n+1$,
so it separates them for every $n$.
With $s=2n+w$, the contour becomes the fixed line
$\mathcal C=\tfrac12+i\mathbb R$.
This is a change of variable; no poles are crossed.
The integral converges absolutely for $z>0$.

Apply Euler's beta integral \cite[Eq.~5.12.1]{NIST:DLMF} to
\[
\frac{\Gamma(w+2n)}{\Gamma(w+13n+1)},\qquad
\frac{\Gamma(w+3n)}{\Gamma(w+14n+1)},\qquad
\frac{\Gamma(w+4n)}{\Gamma(w+12n+2)}.
\]
This gives
\[
\begin{aligned}
\mathcal{G}_n(z)
&=\frac{z^{-2n}}{2\pi i(11n)!^2(8n+1)!}
\int_{\mathcal C}
\Gamma(w)\Gamma(w+n)\Gamma(14n+1-w)\Gamma(15n+1-w)z^{-w}\\
&\quad\times\left[\int_{[0,1]^3}
x^{w+2n-1}y^{w+3n-1}t^{w+4n-1}
(1-x)^{11n}(1-y)^{11n}(1-t)^{8n+1}
\,dx\,dy\,dt\right]dw.
\end{aligned}
\]
All beta parameters have positive real parts on $\mathcal C$.
Apply Euler's integral once more:
\[
\Gamma(w+n)\Gamma(15n+1-w)
=(16n)!\int_0^1 r^{w+n-1}(1-r)^{15n-w}\,dr.
\]
Absolute convergence permits interchanging the integrals.
Writing $u=xyt$, the remaining integral on the same contour is
\cite[Eq.~5.13.1]{NIST:DLMF}
\[
\frac{1}{2\pi i}\int_{\mathcal C}
\Gamma(w)\Gamma(14n+1-w)
\left(\frac{z(1-r)}{ur}\right)^{-w}dw
=(14n)!\left(\frac{ur}{z(1-r)+ur}\right)^{14n+1}.
\]
Its argument is positive and its contour condition
$0<\tfrac12<14n+1$ holds for every $n$.
Consequently, we obtain the positive integral
\[
\mathcal{G}_n(z)=z^{-2n}C_n\int_{[0,1]^4}
\frac{
x^{16n}y^{17n}t^{18n}r^{15n}
(1-x)^{11n}(1-y)^{11n}(1-t)^{8n+1}(1-r)^{15n}
}{[z(1-r)+xytr]^{14n+1}}
\,dx\,dy\,dt\,dr.
\]
In particular, $z=1$ can be substituted directly.
Since $1-r+xytr\geq xyt$, for $n\geq1$ Euler's integral gives
\[
\begin{aligned}
0<\mathcal{G}_n(1)
&\leq C_n\int_{[0,1]^4}
x^{2n-1}y^{3n-1}t^{4n-1}r^{15n}
(1-x)^{11n}(1-y)^{11n}(1-t)^{8n+1}(1-r)^{15n}
\,dx\,dy\,dt\,dr\\
&=\frac{(16n)!(15n)!^2(2n-1)!(3n-1)!(4n-1)!}
{(13n)!(12n+1)!(30n+1)!}
=:D_n.
\end{aligned}
\]
The factorial ratios satisfy, for $n\geq2$,
\[
\frac{D_{n+1}}{D_n}
\leq
\frac{16^{16}15^{30}2^2 3^3 4^4}{13^{13}12^{12}30^{30}}
\left(\frac{n+1}{n}\right)^{55}
\leq
\frac{16^{16}15^{30}2^2 3^3 4^4}{13^{13}12^{12}30^{30}}
\left(\frac32\right)^{55}
=:\eta\approx8.517\cdot10^{-4}<1.
\]
Here the $55$ numerator factors are bounded above using $n+1$,
and the $55$ denominator factors are bounded below using $n$.
Iterating gives, for a constant $C>0$ independent of $n$,
\[
0<\mathcal{G}_n(1)\leq D_n\leq C\eta^n\longrightarrow0.
\]

Finally, $\theta^k$ inserts the factor $(-s)^k$ in the Mellin--Barnes
integrand. For $k\leq3$, these integrals
remain absolutely convergent at $z=1$ on the same contour.
Put $a=2n$. After $s=a+w$, the factor is
$(-1)^k(w+a)^k$.
For $n\geq1$ and $k\leq3$, the polynomial $(w+a)^k$ has an expansion
in the rising factorials $(w)_j$ with nonnegative coefficients.
Each $(w)_j$ replaces $\Gamma(w)$ by $\Gamma(w+j)$; the same Mellin
inversion then multiplies the positive integrand at $z=1$ by
$(L)_j[(1-r)/(1-r+ur)]^j\leq(L)_j$.
Evaluating the polynomial expansion at $w=L$ bounds the combined
factor by $(a+L)^k$. Hence
\[
|(\theta^k\mathcal{G}_n)(1)|
\leq
(16n+1)^kD_n
\leq C(16n+1)^k\eta^n\longrightarrow0,
\qquad 0\leq k\leq3.
\]
Since $16n+1=O(n)$, the claim follows.
\end{proof}

\begin{proof}[Proof of Theorem~\ref{thm:z34-limit}]
By Proposition~\ref{prop:z34-growth}, every component of the
state vector of $(G^{5,2}_{5,5}(1),y_0,v_G)_F$ decays exponentially, so by
Proposition~\ref{prop:z34-rates} it does not grow at any of the rates
$|\rho_1|,|\rho_2|,|\rho_3|$; hence the state is the recessive solution.
Write $\widehat{\mathcal S}_G(n)$ for its first three components and
$\widehat I_G$ for the first three columns of $I_G$.
The state relation and the block form in
Proposition~\ref{prop:z34-coboundary} give
\[
\mathcal S_G(n+1)=\mathcal S_G(n)G_{y_0,v_G}(n),
\qquad
\widehat{\mathcal S}_G(n+1)=\widehat{\mathcal S}_G(n)Q(n).
\]
Hence $\widehat{\mathcal S}_G(n)$ is the recessive solution of $Q(n)$
obtained from the actual Meijer G state. Since $I_G=(\widehat I_G\;0)$,
\[
\widehat I_G\widehat{\mathcal S}_G(0)^T
=I_G\mathcal S_G(0)^T
=6\begin{pmatrix}A\zeta(3)+B\zeta(4)\\1\end{pmatrix}.
\]
The block form also gives the exact restriction
\[
I_GW_G(N)^{-T}
=\bigl(\widehat I_GW_Q(N)^{-T}\;0\bigr),
\qquad W_Q(N)=Q(0)\cdots Q(N-1).
\]
Thus the initial conditions already exclude the $-1$ mode; the
rank-three connection transfers this restriction to the forward walk.
By Theorem~\ref{thm:dual-convergence} and
Proposition~\ref{prop:cmf-initial-limit}, the dual scheme restricted to
$Q(n)$, with initial condition $\widehat I_G$, converges to
$(I_G\mathcal{S}_G(0)^T)_1/(I_G\mathcal{S}_G(0)^T)_2$.
Proposition~\ref{prop:z34-initial-values} proves that this ratio is
$A\zeta(3)+B\zeta(4)$.

At the source point, the coboundary of
Proposition~\ref{prop:z34-coboundary} satisfies
\[
E(0)=\begin{pmatrix}1&0&0\\0&0&0\\0&1&0\\0&0&1\end{pmatrix},
\qquad H(0)=\begin{pmatrix}0&0&1\\-1&2&-1\\-2&2&0\end{pmatrix},
\]
\[
I_FE(0)H(0)=2\widehat I_G.
\]
Equivalently, $I_FK(0)=2I_G$.
Consequently,
\[
I_FE(0)H(0)\widehat{\mathcal S}_G(0)^T
=12\begin{pmatrix}A\zeta(3)+B\zeta(4)\\1\end{pmatrix}.
\]
Put $e_0=(1,0,0)^T$ and $D(n)=\prod_{j=8}^{12}(jn+1)$.
The supplied code verifies
\[
H(n)^{-1}e_0=\frac{1}{D(n)}
\begin{pmatrix}1\\10n+1\\(10n+1)^2\end{pmatrix},
\qquad K(n)J^TH(n)^{-1}e_0=v.
\]
Thus the exact walk identity gives
\[
I_FW_F(N)v
=2\left(\prod_{n=0}^{N-1}\frac{A(n)}{B(n)}\right)
\widehat I_G W_Q(N)^{-T}H(N)^{-1}e_0.
\]
The common scalar cancels in the ratio. For
$T_n=\operatorname{diag}(1,n,n^2)$ and
$Q_\infty=\lim_{n\to\infty}T_nQ(n)T_{n+1}^{-1}$, the transported
projection satisfies $D(n)T_n^{-1}H(n)^{-1}e_0\to w=(1,10,100)^T$.
The supplied code verifies exactly that $Q_\infty$ has distinct
eigenvalues and
\[
\det\bigl[w,\ Q_\infty^{-T}w,\ (Q_\infty^{-T})^2w\bigr]\ne0.
\]
Thus $w$ has a nonzero coefficient along every eigenvector in an
eigenbasis of $Q_\infty^{-T}$, including the eigenvector for the unique
dominant dual root $\rho_4^{-1}$. By the asymptotics of the system at
infinity, this contribution survives the exact intertwining identity;
the rational factors change growth only by polynomial factors.
The denominator's pairing with the initial state is nonzero, so
this gives $R_N\to A\zeta(3)+B\zeta(4)$.
Finally, Proposition~\ref{prop:z34-rates} gives the exponential error
factor $\tau=|\rho_4|/|\rho_1|$, since $|\rho_1|=|\rho_2|$.
\end{proof}

